\abstract

\keywords{РАСПРЕДЕЛЁННЫЕ СИСТЕМЫ, МИКРОСЕРВИСНАЯ АРХИТЕКТУРА, ПЛАТФОРМА УВЕДОМЛЕНИЙ, МУЛЬТИКАНАЛЬНАЯ ДОСТАВКА, GRPC, PROTOBUF, APACHE KAFKA, POSTGRESQL, MONGODB, REDIS, CLICKHOUSE, TRANSACTIONAL OUTBOX, ИДЕМПОТЕНТНОСТЬ, ДЕДУПЛИКАЦИЯ, ПОВТОРНЫЕ ПОПЫТКИ, НАБЛЮДАЕМОСТЬ}

Выпускная квалификационная работа посвящена проектированию и реализации распределённой платформы уведомлений, в которой приём команды, канальная доставка и фиксация результата разделены между сервисными контурами. Проблема работы связана с тем, что уведомительное событие должно обрабатываться при частичных отказах, повторной доставке сообщений брокером и неоднородных каналах связи.

Цель работы состоит в том, чтобы спроектировать и реализовать распределённую платформу уведомлений, обеспечивающую сохранность принятого уведомительного события, управляемую многоканальную обработку и прослеживаемость итогового статуса. Гарантированность доставки рассматривается в границах самой платформы и не распространяется на физическую доставку сообщения конечному пользователю.

В работе исследованы требования к платформам уведомлений, выполнено сравнение архитектурных вариантов и выбран подход с разделением gRPC-командного контура и Kafka-контура фоновой обработки. Реализованы сервисы notification-facade, template-registry, profile-consent, delivery-dispatcher, cancellation-service, notification-mail-sender, notification-sms-sender и history-writer.

Практическим результатом работы является программная платформа, поддерживающая два канальных контура: email и SMS. Реализация включает приём уведомительных событий, мгновенный и отложенный запуск доставки, рендеринг шаблонов, централизованную проверку профиля получателя, повторные попытки при временных ошибках, fallback-маршрутизацию и запись итоговых статусов в историю.

Проверка включала автоматизированные тесты и нагрузочный прогон полного маршрута обработки. Были проверены линейная нагрузка мгновенной отправки, постоянная нагрузка и отложенная доставка. Полученные результаты показали прохождение маршрута от входной команды до записи статуса доставки.
